% \numberwithin{equation}{section}

\chapter{Introducció} \label{sec:intro}

Remember Moncho's advice of using setnewcommand or newcommand to define new commands and make the writing faster.

EVERYTHING HERE IS AI GENERATED, NOT EVEN WITH A PROMPT BUT WITH THE AUTOCOMPLETE FEATURE OF GEMINI COPILOT IN VSCODE. REWRITE IT ALL IN YOUR OWN WORDS.

\section{Spectral distribution of the CMB}
The CMB is the afterglow of the Big Bang, a relic radiation that fills the universe and provides a snapshot of the early cosmos.
It is a nearly perfect blackbody spectrum, with a temperature of approximately 2.725 K.
The spectral distribution of the CMB can be described by Planck's law, which gives the intensity of radiation as a function of frequency or wavelength.

The current upper limits on $\mu$-distortions from COBE-FIRAS are of the order of $|\mu|~<~9~\times~10^{-5}$, which means that any deviations from a perfect blackbody spectrum must be very small.
However, the $\Lambda$CDM model predicts that $\mu$-distortions should be present at a level of around~$10^{-8}$, which is well below the current observational limits.
Detecting $\mu$-distortions would provide valuable insights into the physical processes that occurred in the early universe,
such as energy release from decaying particles, dissipation of acoustic waves, or other non-standard processes that could have injected energy into the photon-baryon plasma before recombination



\section{SZ effect}
The Sunyaev-Zeldovich (SZ) effect is a distortion of the CMB spectrum caused by the inverse Compton scattering of CMB photons off high-energy electrons in galaxy clusters.
This effect can be used to study the properties of galaxy clusters and the large-scale structure of the universe.
The thermal SZ effect, in particular, is sensitive to the pressure of the intracluster medium and can be used to measure the cluster's gas content and temperature.
Since the SZ effect involves the scattering of CMB photons, any primordial $\mu$-distortion would also imprint a corresponding $\mu$-dependent modification in the SZ signal.
This means that if we can detect the SZ effect from galaxy clusters, we can also look for the presence of $\mu$-distortions in the SZ signal,
which could provide a new way to constrain $\mu$-distortions and gain insights into the early universe.



\section{$\mu$-distorted SZ effect}
The $\mu$-distorted SZ effect refers to the modification of the thermal SZ effect due to the presence of $\mu$-type distortions in the CMB spectrum.
If the CMB photons that are scattered by the hot electrons in galaxy clusters have a $\mu$-distortion, then the resulting SZ signal will also be modified in a way that depends on the value of $\mu$.
This means that by analyzing the SZ signal from galaxy clusters, we can potentially detect or constrain the presence of $\mu$-distortions in the CMB.


% Chain... if SZ effect is there, then $\mu$-distorted SZ effect is there too.

\section{Other effects}
There are many other effects that can modify the CMB spectrum, such as y-distortions from the thermal SZ effect, kinematic SZ effect, and foreground emissions from our galaxy and extragalactic sources \cite[Zegeye2024_mu_distortion].
These effects can complicate the detection of $\mu$-distortions, as they can mimic or mask the signal we are trying to detect.
Therefore, it is important to use component separation techniques, such as the Constrained Internal Linear Combination (CILC), to isolate the $\mu$-distorted SZ signal from these other contributions.
By applying these techniques to sky simulations for future CMB surveys, we can assess whether this SZ-based approach can improve current upper limits on $\mu$-distortions and potentially lead to a detection of these distortions in the future.
